\section{Discussion}
\label{sec:discussion}

The experimental results show that synchronized BiLSTM trajectory, ViT visual, and AST audio signals provide complementary information for \task{}. This supports the central assumption of the study: fine-grained combine-harvester states cannot be reliably described by spatial-temporal motion alone because visually and mechanically different operations may share similar rows, speeds, stops, passes, and turns. Direct trimodal concatenation improves the fixed-seed macro-F1 by 6.56 percentage points over BiLSTM and by 14.62 points over AST, indicating that the gain is not only an accuracy effect from frequent classes but also a more balanced improvement across the eleven-class taxonomy. The proposed class-adaptive gate further improves macro-F1 by 2.44 points over concatenation, reaching 63.64\% macro-F1 and 81.36\% accuracy.

The per-class results further explain why each modality is useful. BiLSTM provides spatial continuity, direction, turning behavior, and transfer movement, which are essential for reverse transfer and straight transfer. ViT provides crop contact, road surface, visible harvesting conditions, and field-road context, which benefits empty harvesting and road driving. AST provides engine response, load variation, and working-component engagement, which explains its strong performance on full-load harvesting, engine-off waiting, idling waiting, and unloading. These findings are consistent with the practical harvesting workflow and show that multimodal sensing is not only a model-level improvement but also an operationally interpretable extension of trajectory-only classification.

At the same time, the fusion ablation reveals both the value and the limitation of class-adaptive gating. Compared with concatenation, the gate improves turning transfer by 11.74 F1 points and idling waiting by 9.60 points, and it reduces the dominant idling waiting $\rightarrow$ engine-off waiting error by 791 samples. This supports the design choice of keeping modality-specific logits before the final decision. However, unloading and road driving remain below their best single-modality baselines. These classes are either strongly process-dominated or rare in the current split, so the learned gate can still dilute the most reliable modality. This motivates future constraints on class-wise gate weights or prior-guided fusion for rare and audio-dominant states.

The dataset also reflects naturally imbalanced agricultural production rather than an artificially balanced benchmark. Full-load harvesting and waiting occupy large temporal proportions, whereas turning transfer and road driving are minority classes. This imbalance explains why accuracy alone can be misleading: AST reaches 74.24\% accuracy but only 46.57\% macro-F1 because it performs well on major process classes while failing on several minority motion classes. Macro-F1, per-class recall, and trajectory-level diagnostics are therefore necessary for evaluating the practical usefulness of \task{} models.

Several limitations remain. First, the final trimodal comparison currently uses a fixed seed, while only the BiLSTM sequence-length ablation has been evaluated across multiple seeds. Additional final-suite seeds are needed before making strong statistical claims. Second, the study is restricted to combine harvesters during maize harvest operations in three regions; cross-machine, cross-crop, and cross-season generalization remains to be tested. Third, the current multimodal model uses synchronized observations but does not explicitly model transition boundaries or class-neighborhood structure, which contributes to errors such as turning transfer being confused with reverse or straight transfer. Future work should therefore combine the class-adaptive gate with transition-aware temporal modeling and imbalance-aware objectives for minority operation classes.
